% Part: proof-theory
% Chapter: sequent-calculus
% Section: interpretation-rules

\documentclass[../../../include/open-logic-section]{subfiles}

\begin{document}

\olfileid{pt}{seq}{irl}

\olsection{Interpretasi Aturan}

Elinga interpretasi sekuen: sekuen $\Gamma \Sequent \Delta$ bener
ing !!a{structure}~$\Struct{M}$ yen paling ora ana siji $!A \in \Gamma$
sing salah utawa paling ora ana siji~$!A \in \Delta$ sing bener.
Aturan logis~\Log{G3c} bisa diwaca minangka rumusan syarat kabeneran
kanggo !!{formula}s utamane. Gatekna \RightR{\lif}. !!{formula}
utamane yaiku $!A \lif !B$. Formula iki bener yen $!A$ salah
utawa~$!B$ bener. Mula sekuen $\ \Sequent !A \lif !B$ bener yen
lan mung yen $!A \Sequent !B$ bener. Yen !!{formula}s konteks
$\Gamma$, $\Delta$ ditambahake ing anteseden lan sukseden sekuen
iki, kita oleh persis kesimpulan lan premis aturan \RightR{\lif}.
Kahanane padha kanggo~\LeftR{\lif}. $!A \lif !B$ salah yen
$!A$~bener lan $!B$~salah. Dadi $!A \lif !B \Sequent \ $ bener
yen lan mung yen $\ \Sequent !A$ lan $!B \Sequent \ $ loro-lorone
bener. Kesimpulan \LeftR{\lif} ekuivalen karo konjungsi
premis-premise. Pola iki njamin aturan~\Log{G3c} sahih: yen kabeh
premis bener, kesimpulane bener. Pola iki uga njamin kelengkapan.

Satemene, syarat kabeneran saben panyambung fungsional-kabeneran
bisa diterangake kanthi cara iki minangka kumpulan sekuen.
Iki amarga saben fungsi kabeneran ing logika klasik bisa diterangake
kanthi !!a{formula} awujud normal konjungtif: konjungsi disjungsi
saka !!{formula}s atomik lan negasi formula atomik. Upamane
$!A \oplus !B$ bener yen lan mung yen persis siji saka $!A$ lan
$!B$ bener, yaiku ``utawa eksklusif.'' Mula $!A \oplus !B$ bener
yen lan mung yen $(!A \lor !B) \land (\lnot !A \lor \lnot !B)$
bener, lan salah yen $(\lnot !A \lor !B) \land (!A \lor \lnot !B)$
bener. Dadi kita bisa ngenalake aturan kalkulus sekuen kanggo
$\oplus$ kaya iki:
\begin{align*}
&
\Axiom$\Gamma \fCenter \Delta, !A, !B$
\Axiom$!A, !B, \Gamma \fCenter \Delta$
\RightLabel{\RightR{\oplus}}
\BinaryInf$\Gamma \fCenter \Delta, !A \oplus !B$
\DisplayProof
\\
& \Axiom$!A, \Gamma \fCenter \Delta, !B$
\Axiom$!B, \Gamma \fCenter \Delta, !A$
\RightLabel{\LeftR{\oplus}}
\BinaryInf$!A \oplus !B, \Gamma \fCenter \Delta$
\DisplayProof
\end{align*}
Aturan iki uga sahih lan jangkep. Kesimpulane bener yen lan mung yen
kabeh premise bener.

\begin{prob}
Upamane $!A \downarrow !B$ bener yen lan mung yen $!A$ lan~$!B$
loro-lorone ora bener (``NOR''). Kepiye aturan \Log{G3c} kanggo
panyambung iki? (Golekana rong sekuen sing bener bebarengan yen lan
mung yen $!A \downarrow !B$ bener. Iki menehi aturan
\RightR{\downarrow}. Banjur golekana siji sekuen sing bener yen
lan mung yen $!A \downarrow !B$ salah. Iki menehi aturan
\LeftR{\downarrow}.)
\end{prob}

Ing \Log{G3c}, saben panyambung lan kuantor duwe persis rong aturan:
aturan kiwa lan aturan tengen. Nanging ing \Log{G1c}, ana rong aturan
\LeftR{\land} lan rong aturan \RightR{\lor}. Sebabe yaiku sistem
asli Gentzen~\Log{LK}, sing dadi tuladha kanggo \Log{G1c}, duwe
varian~\Log{LJ} kanggo logika nonklasik sing wigati, yaiku logika
intuisionistik. Supaya sistem sahih lan jangkep tumrap logika
intuisionistik, sukseden saben sekuen ing~\Log{LJ} diwatesi supaya
ngemot paling akeh siji !!{formula}. Mula aturan \RightR{\lor}
saka~\Log{G3c} ora bisa digunakake, amarga premise ngemot $!A$ lan~$!B$
ing sukseden. Dadi Gentzen nggunakake sapasang aturan \RightR{\lor}
kanggo \Log{LJ}, lan pasangan sing padha kanggo~\Log{LK}.
Semono uga, versi intuisionistik~\Log{G1i} yaiku~\Log{G1c} kanthi
kabeh sukseden diwatesi supaya ngemot paling akeh siji !!{formula}.
(\Log{G3c} uga duwe versi intuisionistik, nanging sawetara aturane
beda saka aturan sing cocog ing~\Log{G3c}. Upamane, versi iku uga
nggunakake sapasang aturan \RightR{\lor}.)

Gampang dimangerteni yen aturan struktural~\Log{G1c} sahih.
Upamane, yen $\Gamma \Sequent \Delta$ ngemot !!{formula} sing salah
ing sisih kiwa utawa !!{formula} sing bener ing sisih tengen, mula
$\Gamma \Sequent \Delta, !A$ uga ngemot formula sing padha; apa
$!A$ bener utawa salah ora mengaruhi iki. Dadi yen premis
aturan~\RightR{\Weakening} bener, kesimpulane uga bener. Gatekna
yen arah walikane ora bener: kesimpulan bisa bener amarga $!A$
bener, lan ing kahanan iki premise ora mesthi bener.

Ngapa aturan struktural dibutuhake? Kanggo formula skematis atomik saka predikat
sing beda-beda, kontraksi (\RightR{\Contraction}, \LeftR{\Contraction})
dibutuhake, upamane, kanggo !!{prove} $\ \Sequent !A \lor \lnot !A$.
Yen diwiwiti saka $!A \Sequent !A$, dhisik kita oleh $\Sequent !A,
\lnot !A$ kanthi~\RightR{\lnot}. Banjur rong versi \RightR{\lor}
bisa digunakake runtut: sapisan kanggo oleh $!A \lor \lnot !A$ saka
$\lnot !A$, lan sapisan kanggo oleh $!A \lor \lnot !A$ saka~$!A$.
Asile yaiku $\ \Sequent !A \lor \lnot !A, !A \lor \lnot !A$.
Dadi supaya oleh $\Sequent !A \lor \lnot !A$ ing~\Log{G1c}, kita
kudu bisa ``ngontraksi'' rong kedadeyan !!{formula} sing padha
ing !!a{proof}:
\[
\Axiom$!A \fCenter !A$
\RightLabel{\RightR{\lnot}}
\UnaryInf$\fCenter !A, \lnot !A$
\RightLabel{\RightR{\lor}}
\UnaryInf$\fCenter !A, !A \lor \lnot !A$
\RightLabel{\RightR{\lor}}
\UnaryInf$\fCenter !A \lor \lnot !A, !A \lor \lnot !A$
\RightLabel{\RightR{\Contraction}}
\UnaryInf$\fCenter !A \lor \lnot !A$
\DisplayProof
\]
Pancen ana sekuen valid sing ora bisa !!{prove} ing~\Log{G1c} tanpa
pelemahan utawa kontraksi. Upamane, !!a{proof} kanggo sekuen
$!A \Sequent !B \lif !A$ ing~\Log{G1c} kudu dipungkasi nganggo
\RightR{\lif}, amarga mung $!B \lif !A$ sing bisa dadi formula utama,
lan aturan iki mbutuhake $!B, !A \Sequent !A$ minangka premis.
Premis iki mbutuhake \LeftR{\Weakening} supaya bisa dijupuk saka
aksioma $!A \Sequent !A$:
\[
\Axiom$!A \fCenter !A$
\RightLabel{\LeftR{\Weakening}}
\UnaryInf$!B, !A \fCenter !A$
\RightLabel{\RightR{\lif}}
\UnaryInf$!A \fCenter !B \lif !A$
\DisplayProof
\]

Ing \Log{G3c}, aturan kontraksi lan pelemahan ora dibutuhake
babar pisan. Pelemahan wis kalebu ing aksioma. Kontraksi umume
bisa disingkiri kanthi nggabungake rong salinan konteks $\Gamma$,
$\Delta$ ing aturan rong premis dadi siji. \Log{G1c} lan \Log{G3c}
loro-lorone duwe aturan sing ``nuduhake konteks sing padha'' iki.
Ing aturan siji premis, kontraksi mung dibutuhake kanggo
\RightR{\lor}, \LeftR{\land}, \RightR{\lexists}, lan~\LeftR{\lforall}.
Versi aturan iki ing~\Log{G3c} ndadekake kontraksi ora perlu kabeh.
Ana sistem \Log{G2c} (delengen \olref{tab:G2c}) sing aturan rong
premise duwe konteks mandhiri. Ing \Log{G2c}, kontraksi luwih wigati
tinimbang ing~\Log{G1c}.

\subfile{rules-G2c}

\end{document}
